\section{The 195-state construction and its preservation theorem}
\label{sec:preservation195}

Let $M_{196}^{\mathrm{sep}}$ be the separated-diagonal machine supplied in
\texttt{pa\_consistency\_comparison\_and\_new196\_2026-11-13.zip},
with the original PA196 command control. Let $M_{195}$ be the new table
in this archive using the same fixed control and the new axiom-reader
grammar. Both machines use a two-way infinite binary tape, initially all
zero. In native labels the initial state is $74$ and the unique halting
label is $254$. A dense relabeling does not affect the theorem.

\paragraph{What the construction changes.}
The new grammar has 67 distinct reader nodes. The fixed part has 134
working states, and six disjoint, simultaneously valid context folds
remove six reader states. Hence the number of working states is
\[
 134+67-6=195.
\]
There are exactly 390 transition rows, one for each working-state and
symbol pair. Every working state is graph-reachable from the initial
state. The halting label is not counted as a working state.

The flattened axiom-emitter word has exactly the same 190 tokens, in
the same order, as the supplied separated-diagonal emitter. Although the
search allowed independent parameter permutations, the final grammar
uses the identity permutation in all six heads. No axiom formula,
parameter assignment, guard, stack position, or command selector changes.
The new object is a different binary parenthesization and sharing of the
same emitter. All fixed transition rows are unchanged except the reader
entry transition $(101,1)$, which is redirected to the new reader root.

Token equality alone would be insufficient for this tree language:
changing ancestors can change the targets of structured jumps. The
stronger check is exact equality of the entire branch-resolved register
program. Both source trees have 3,673 primitive sites. Number their leaves
from left to right by $0,\ldots,3672$. At every ordinal, the instruction,
register operand, and every successor ordinal agree. This includes both
branches of every decrement, every break and continue target, and every
halt. Only source addresses in the binary tree change. The equality is
checked using both arithmetic path cutting and an independent resolver
based on subtree intervals and ancestor walking. The latter also checks
that a proposed successor has no nonzero suffix beyond its target leaf.

\paragraph{Source identity lemma.}
Let $G$ be this common finite control graph, acting on nine natural-number
registers. The two source programs have identical executions after
identifying their program counters by leaf ordinal. This holds for every
initial register vector, not merely the clean seed. Indeed, the step
functions are literally the same finite case distinction: increment has
the same register and successor; decrement tests and updates the same
register and has the same two successors; a jump has the same target;
and halt has the same meaning. Induction on the number of source steps
proves equality of their finite traces. Since the graph is finite and
explicit, PA proves the same statement about its coded step functions.

\paragraph{The inherited tape interfaces.}
The source-to-tape argument uses the original PA196 interfaces for
scanning and marking a unary register, incrementing it, and decrementing
it with its zero/positive case distinction. These interfaces quantify
over the complete arbitrary register file and the intervening tape gap.
They specify the action on that whole file: in particular, increment
also appends a one-valued spare register, as in the inherited engine.
The untouched logical registers are preserved, and the file continues
to satisfy its minimum-length invariant. The interfaces also assert
termination. The register routines in $M_{195}$ are literally unchanged.
More precisely, the new table agrees on all 44 entries of
\texttt{M6.Native.boxRows} and on the append-one tail
$(198,0)\mapsto(1,L,40)$. The parameterized theorem
\texttt{M6.Native.boxLemmas} therefore supplies the same four register
contracts for this table. In the upstream formal development these are
the block contracts supplied to the symbolic executor in
\texttt{M6/Runtime.lean}, whose soundness chain includes
\texttt{oracle\_sound}, \texttt{sstep\_sound}, and
\texttt{exec\_sound}. The present executable audit uses those mathematical
block interfaces; it does not replace them by bounded tests on numerical
register samples.

For each primitive source site, and for both branches when appropriate,
the delivered symbolic certificate begins at the full encoded source
frame. It reaches exactly the next full frame with the specified register
effect, or reaches the halting state for a halt instruction. The symbolic
gap is affine in an arbitrary natural number, and the register file is
an arbitrary symbolic object satisfying the stated branch condition.
The 195-state certificate has 4,780 cases, a maximum of 452 symbolic
steps in a case, and 1,442,472 symbolic steps in total. These are symbolic
proof-case counts, not bounds on the physical execution time for
arbitrarily large register values.

Every nonhalting source instruction is simulated by a finite, positive
number of physical steps, without a premature halt. The positive-length
condition is explicit in the executor: reaching the target does not
close a case until an actual tape transition has occurred. Compressed
constant-symbol scans and register oracles represent finite runs whose
lengths have elementary, hence primitive recursive, bounds in the frame
parameters. This excludes an infinite physical computation trapped
inside one source instruction. The complete source graph audit supplies
the correct target for each symbolic case.

The six reader folds have an additional local check. Before folding, each
member of a pair is entered by a rightward, symbol-preserving move, with
the bit just crossed distinguishing the two members. One member has only
incoming context zero, and the other only incoming context one. For each
of these two contexts and each current symbol, the replacement follows
exactly the original transition's tape effect, head displacement and
target, allowing a finite detour that inspects the previous bit. All 24
cases are checked against the final table after all six folds. The
detours stay within the two cells described by the local interface, so
the arbitrary tape outside them is untouched. The whole-frame symbolic
certificate then checks the folded table itself, not only its unfolded
parent.

\paragraph{Blank initialization.}
Both the supplied parent and the new table reach their exact clean source
frames after 58,154 literal tape transitions from the all-zero tape.
The source counter is site 25 and the logical register vector is
$(2,0,0,0,0,0,0,0,0)$. The complete tape, head position and native dispatch
state 84 are checked, including the unchanged dummy-register suffix.
This is a fixed finite trace, so its correctness is an elementary PA
calculation. Agreement only on a selected portion of the tape would
not suffice here; the delivered check compares the complete finite
support of the tape.

\paragraph{Finite-trace equivalence theorem.}
Write $H_M$ for the arithmetic sentence that $M$ has a finite halting
trace from its specified blank start. Then
\[
 \mathrm{PA}\vdash
 H_{M_{195}}\ \longleftrightarrow\ H_{M_{196}^{\mathrm{sep}}}.
\tag{195-preservation}
\]
To prove this inside PA, use the common source graph $G$ and the exact
clean seed. A finite halting source trace produces a finite halting tape
trace by concatenating the exact startup and primitive execution
segments. Their codes and finite bounds are primitive recursively
constructible, and PA proves their correctness by induction on the
length of the source trace.

For the reverse implication from a tape trace to a source trace, suppose
the tape trace first halts after $T$ transitions. Iterate the common
source step function for $T+1$ instructions from the clean seed. If none
of these instructions halts, their exact positive-length simulations,
after startup, concatenate to a physical run longer than $T$, with no
halt before its end. Determinism contradicts the alleged first halting
trace. Thus one of the first $T+1$ source instructions halts. This
bounded iteration is primitive recursive; its correctness needs only
ordinary arithmetic induction. Apply both directions to each of the
two machines and eliminate the common source-halting predicate. The
argument does not assume that either machine actually fails to halt.

\paragraph{Composition with the PA calibration.}
The numerical calculus and proof-search program in this theorem are
exactly the supplied separated-diagonal ones. The inherited mathematical
calibration gives
\[
 \mathrm{PA}\vdash
 H_{M_{196}^{\mathrm{sep}}}
 \ \longleftrightarrow\
 \exists p\,\operatorname{Prf}_{\mathrm{PA}}
       (p,\ulcorner 0=1\urcorner).
\tag{parent-calibration}
\]
Its substantive interfaces are the following. The arithmetic
interpretation translates membership to finite binary-set membership
and supplies PA proofs of pointed separation by full induction. The
formula-code freshness lemma supplies proofs of the diagonal clause
$B\to U(B,B)$. The recovered two-direction membership rewrites,
neighborhood equality and equality-witness interfaces feed the original
PA-refutation compiler. The separated diagonal head uses the adjusted
ten-item stack selectors. In the soundness direction, a proof annotation
of the top, cached head and canonical Cantor stack turns a source halt
into a PA refutation, including the zero-stack alias. In the completeness
direction, the proof-command compiler and exhaustive command scheduler
turn a PA refutation into a source halt. These are primitive recursive
syntactic constructions with PA-provable correctness; no uniform truth
predicate or reflection principle for PA is used. The inherited
calibration and diagonal-interface arguments are included elsewhere in
the report so that the logical dependency is visible.

Combining parent calibration with (195-preservation) and the definition
$\operatorname{Con}(\mathrm{PA})\equiv
\neg\exists p\,\operatorname{Prf}_{\mathrm{PA}}
(p,\ulcorner0=1\urcorner)$ gives the requested result:
\[
 \boxed{\displaystyle
 \mathrm{PA}\vdash
 \bigl(M_{195}\text{ does not halt from its blank start}\bigr)
 \ \longleftrightarrow\ \operatorname{Con}(\mathrm{PA}).}
\]

This is an ordinary mathematical proof obtained by preservation and
composition with the stated upstream interfaces. The executable files
make the finite construction and its certificate data reproducible.
They do not constitute a newly exported Lean theorem for the entire
displayed biconditional. The inherited formal development and the
additional written arithmetic-calibration argument are distinct parts
of the proof provenance.
